\section{The Four Endowments of the Glorified Body}

What manner of flesh enters the Father's house? St.\ Paul sets our mortal frame beside the risen body in four stark antitheses:

\begin{quote}
It is sown in corruption: it shall rise in incorruption. It is sown in dishonour: it shall rise in glory. It is sown in weakness: it shall rise in power. It is sown a natural body: it shall rise a spiritual body. (1 Cor 15:42--44)
\end{quote}

From these four contrasts and the Easter appearances of Jesus Christ---whose risen humanity is the cause and exemplar of ours (\textit{ST} III, q.\ 54, a.\ 2; Phil 3:21)---the Fathers and Scholastic Doctors draw out the four endowments (\textit{dotes}) of the glorified body: \textbf{impassibility}, \textbf{clarity}, \textbf{agility}, and \textbf{subtlety}.\endnote{\textit{Catechism of the Council of Trent}, Part I, Creed Article XI, trans.\ J.\ Donovan (Baltimore: Lucas Brothers, 1829), pp.\ 90--92; cf.\ St.\ Thomas Aquinas, \textit{Summa Theologiae} III, q.\ 54, aa.\ 1--2.} None of the four is bolted onto matter from outside. As St.\ Thomas explains of Christ, ``Christ's soul was glorified from the instant of His conception by perfect fruition of the Godhead,'' and after the Passion ``straightway the soul communicated its glory to the risen body in the Resurrection; and so that body was made glorious.''\endnote{St.\ Thomas Aquinas, \textit{Summa Theologiae} III, q.\ 54, a.\ 2.} Beatitude belongs first to the soul beholding God, and overflows from the soul into every fiber of the reunited body.

\subsection*{Impassibility}

The \textit{Roman Catechism} begins with \textbf{impassibility} (\textit{impassibilitas}), answering to ``it is sown in corruption: it shall rise in incorruption'' (1 Cor 15:42):

\begin{quote}
The first is ``impassibility,'' which shall place them beyond the reach of pain or inconvenience of any sort. Neither the piercing severity of cold, nor the glowing intensity of heat can affect them, nor can the impetuosity of waters hurt them. ``It is sown,'' says the Apostle, ``in corruption, it shall rise in incorruption.'' This quality, the schoolmen call impassibility, not incorruption: in order to distinguish it as a property peculiar to a glorified body. The bodies of the damned, though incorruptible, shall not be impassible: they shall be capable of experiencing heat and cold, and of feeling pain.\endnote{\textit{Catechism of the Council of Trent}, Part I, Creed Article XI, trans.\ Donovan, pp.\ 90--91.}
\end{quote}

Every dead person rises incorruptible, yet only the blessed receive impassibility: no hunger, thirst, scorching heat, sickness, aging, or death can touch them again (Apoc 7:16; 21:4; Luke 20:36). Impassibility is not stony numbness: as the Fathers and the \textit{Roman Catechism} make clear, the glorified body remains true living flesh whose senses are shielded from injury, suffering, and decay rather than stripped of life (cf.\ \textit{ST} III, q.\ 54, aa.\ 1--2; \textit{Supplement}, q.\ 91, a.\ 1). As St.\ Cyril of Jerusalem writes: ``This body therefore shall be raised, but it shall abide not such as it now is, but an eternal body; no longer needing for its life such nourishment as now, nor stairs for its ascent, for it shall be made spiritual, a marvellous thing, such as we cannot worthily speak of.''\endnote{St.\ Cyril of Jerusalem, \textit{Catechetical Lectures} 18.18, trans.\ Gifford, NPNF2 vol.\ 7.}

\subsection*{Clarity}

The second gift is \textbf{clarity} or \textbf{brightness} (\textit{claritas}), answering to ``it is sown in dishonour: it shall rise in glory'' (1 Cor 15:43). Christ states its measure plainly: ``Then shall the just shine as the sun, in the kingdom of their Father'' (Matt 13:43)---a splendor prefigured on Sinai in the glowing face of Moses (Exod 34:29) and unveiled on Mount Tabor, where ``his face did shine as the sun: and his garments became white as snow'' (Matt 17:2). The \textit{Roman Catechism} traces this radiance to its inward spring:

\begin{quote}
This brightness is a sort of refulgence reflected from the supreme happiness of the soul---an emanation of the bliss which it enjoys, and which beams through the body. Its communication is analogous to the manner in which the soul itself is rendered happy, by a participation of the happiness of God. Unlike the former, this quality is not common to all in the same degree. All the bodies of the saints shall, it is true, be equally impassible: but the brightness of all shall not be the same: for, according to the Apostle: ``One is the glory of the sun, another the glory of the moon, and another the glory of the stars, for star differeth from star in glory: so also, is the resurrection of the dead.''\endnote{\textit{Catechism of the Council of Trent}, Part I, Creed Article XI, trans.\ Donovan, p.\ 91.}
\end{quote}

Does such radiance blot out human features? Answering this question of Christ's risen flesh (\textit{ST} III, q.\ 54, a.\ 2, ad 1), Aquinas compares glorified bodies to colored glass lit by the sun: ``the brightness or splendor of a glorified body is after the manner of natural color in the human body; just as variously colored glass derives its splendor from the sun's radiance, according to the mode of the color.''\endnote{St.\ Thomas Aquinas, \textit{Summa Theologiae} III, q.\ 54, a.\ 2, ad 1.} Every saint shines in his own proper measure while remaining unmistakably recognizable. St.\ Cyril of Jerusalem adds an unforgettable sign from a summer night: ``God, fore-knowing men's unbelief, has given to little worms in the summer to dart beams of light from their body, that from what is seen, that which is looked for might be believed; for He who gives in part is able to give the whole also, and He who made the worm radiant with light, will much more illuminate a righteous man.''\endnote{St.\ Cyril of Jerusalem, \textit{Catechetical Lectures} 18.18, trans.\ Gifford, NPNF2 vol.\ 7.}

\subsection*{Agility}

The third gift is \textbf{agility} (\textit{agilitas}), answering to ``it is sown in weakness: it shall rise in power'' (1 Cor 15:43). Now ``the corruptible body is a load upon the soul'' (Wis 9:15); by agility, the \textit{Roman Catechism} teaches, ``the body shall be freed from the burden that now presses it down; and shall [acquire] a capability of moving with the utmost facility and celerity, wherever the soul pleases.''\endnote{\textit{Catechism of the Council of Trent}, Part I, Creed Article XI, trans.\ Donovan, p.\ 91 (the 1829 printing reads ``require'' for the Latin \textit{acquiret}).} Fulfilling Wisdom's promise that ``the just shall shine, and shall run to and fro like sparks among the reeds'' (Wis 3:7), the risen body moves swiftly and effortlessly through space, freed from weight and fatigue. St.\ Augustine states its governing law with classic restraint: ``What power of movement such bodies shall possess, I have not the audacity rashly to define, as I have not the ability to conceive. Nevertheless I will say that in any case, both in motion and at rest, they shall be, as in their appearance, seemly; for into that state nothing which is unseemly shall be admitted. One thing is certain, the body shall forthwith be wherever the spirit wills, and the spirit shall will nothing which is unbecoming either to the spirit or to the body.''\endnote{St.\ Augustine, \textit{The City of God} XXII.30, trans.\ Dods.}

\subsection*{Subtlety}

The fourth endowment is \textbf{subtlety} (\textit{subtilitas}), answering to ``It is sown a natural body: it shall rise a spiritual body'' (1 Cor 15:44). ``Spiritual body'' does not mean mist or ether. St.\ Augustine refutes that misreading in \textit{The City of God} XXII.21:

\begin{quote}
The flesh shall then be spiritual, and subject to the spirit, but still flesh, not spirit, as the spirit itself, when subject to the flesh, was fleshly, but still spirit and not flesh. \dots\ And a man is in this life spiritual in such a way, that he is yet carnal with respect to his body, and sees another law in his members warring against the law of his mind; but even in his body he will be spiritual when the same flesh shall have had that resurrection of which these words speak, ``It is sown an animal body, it shall rise a spiritual body.''\endnote{St.\ Augustine, \textit{The City of God} XXII.21, trans.\ Dods.}
\end{quote}

Subtlety does not dissolve matter; as the \textit{Roman Catechism} states, it is ``a quality which subjects the body to the absolute dominion of the soul, and to an entire obedience to her control.''\endnote{\textit{Catechism of the Council of Trent}, Part I, Creed Article XI, trans.\ Donovan, p.\ 91.} On Easter night Christ entered the Upper Room ``the doors being shut'' (John 20:26) and vanished at Emmaus (Luke 24:31), yet bade the disciples touch Him: ``Handle, and see: for a spirit hath not flesh and bones, as you see me to have'' (Luke 24:39). Explaining that Easter miracle in \textit{ST} III, q.\ 54, a.\ 1, ad 1, Aquinas insists that Christ's risen flesh remained true, palpable matter: ``it was not from any property within the body, but by virtue of the Godhead united to it, that this body, although a true one, entered in among the disciples while the doors were shut.''\endnote{St.\ Thomas Aquinas, \textit{Summa Theologiae} III, q.\ 54, a.\ 1, ad 1.}

\subsection*{How the Bodily Eye Sees in Heaven}

Prompted by Job 19:26--27 (``in my flesh I shall see my God''), St.\ Augustine asks in \textit{The City of God} XXII.29 whether glorified eyes of flesh behold God Himself.\endnote{St.\ Augustine, \textit{The City of God} XXII.29, trans.\ Dods.} The saints certainly do not lose the Beatific Vision when they shut their eyes; yet if bodily eyes see God at all, Augustine judges it far more credible that ``God will be known to us and visible to us in the sense that He will be spiritually perceived by each one of us in each one of us, perceived in one another, perceived by each in himself; He will be seen in the new heaven and the new earth, in the whole creation which shall then exist; He will be seen in every body by means of bodies, wherever the eyes of the spiritual body are directed with their penetrating gaze.''

The \textit{Supplement} (q.\ 91, a.\ 1) adopts that resolution cleanly: because God is incorporeal, ``the carnal eye \dots\ will be unable to attain to this vision of the Essence,'' yet ``that it may be fittingly comforted in the vision of God, it will see the Godhead in Its corporeal effects, wherein manifest proofs of the Divine majesty will appear, especially in Christ's flesh, and secondarily in the bodies of the blessed, and afterwards in all other bodies.''\endnote{\textit{Supplement}, q.\ 91, a.\ 1 (English Dominicans, 1922).} With the mind's eye the saint looks straight into the uncreated Trinity; with the eyes of his risen flesh he beholds the radiant face of Christ, His Mother, every glorified friend, and a renewed creation.
